% (User input window)
%
% (Source tag) /killo-golden
%
% 内容以本文既有记号与可审计接口重写，不引入手工编号。

\paragraph{例外谱幂和与低阶迹矩的指数分解：Lucas--Pell 纠正项、符号型递推与有限枚举核验}
令 \(q\ge 2\)，并取 \(U_q=\{0,1\}^q\)；将 \(U_q\) 与 \(\mathcal P([q])\) 通过指示向量识别，并以 \(u\wedge v\) 表示逐坐标与。
定义对称矩阵 \(T^{(q)}\in\RR^{U_q\times U_q}\) 为
\[
T^{(q)}_{u,v}=
\begin{cases}
1,&u\wedge v=0,\\[2pt]
\frac12,&u\wedge v\neq 0,
\end{cases}
\qquad (u,v\in U_q).
\]
由命题 \ref{prop:pom-multiplicity-composition-sq-symmetry-reduction}，\(T^{(q)}\) 在 \(S_q\)-对称子空间上的谱由显式 \((q+1)\times(q+1)\) 矩阵 \(A^{(q)}\) 给出；记该 \((q+1)\) 个单根特征值（例外谱）为 \(\{\lambda_i^{(q)}\}_{i=0}^{q}\)，并对整数 \(m\ge 1\) 定义幂和
\[
S_m(q):=\sum_{i=0}^{q}\big(\lambda_i^{(q)}\big)^m.
\]

\begin{conclusion}[例外谱一次幂和的整系数塌缩：\texorpdfstring{$S_1(q)=2^{q-1}+\frac12F_{q+1}$}{S1(q)=2^{q-1}+1/2 F_{q+1}}]\label{con:xi-terminal-replica-softcore-exceptional-power-sum-m1-integer-collapse}
令 Fibonacci 数列满足 \(F_0=0\)、\(F_1=1\)。
则对一切 \(q\ge 2\) 成立严格闭式
\[
\boxed{\;
S_1(q)=2^{q-1}+\frac12\,F_{q+1}\; }.
\]
因此例外谱总和在 \(\QQ(\sqrt5)\) 中出现的 \(\varphi\)-表示可完全消去，并在 \(\QQ\) 中塌缩为半整数序列。
\end{conclusion}
\begin{proof}
由于 \(S_1(q)=\Tr(A^{(q)})\)，由 \eqref{eq:pom-multiplicity-composition-Aq-def} 直接计算
\[
\Tr(A^{(q)})
=\frac12\sum_{k=0}^{q}\binom{q}{k}+\frac12\sum_{k=0}^{q}\binom{q-k}{k}
=2^{q-1}+\frac12\sum_{k\ge 0}\binom{q-k}{k}.
\]
再用经典恒等式 \(\sum_{k\ge 0}\binom{q-k}{k}=F_{q+1}\) 即得结论（亦与定理 \ref{thm:pom-replica-softcore-exceptional-spectrum-trace} 等价）。证毕。
\end{proof}

\begin{conclusion}[Lucas--Pell 纠正项的统一生成与符号型递推律]\label{con:xi-terminal-replica-softcore-lucas-pell-omega-unified}
令 Lucas 数列满足 \(L_0=2\)、\(L_1=1\)，并令 \(\varphi=\frac{1+\sqrt5}{2}\)。
对任意整数 \(m\ge 1\) 定义
\[
\Omega_m(q):=
\begin{cases}
\dfrac{\varphi^{m(q+2)}+(-1)^q\varphi^{-mq}}{\varphi^{2m}+1},& m\ \text{为奇数},\\[10pt]
\dfrac{\varphi^{m(q+2)}-\varphi^{-mq}}{\varphi^{2m}-1},& m\ \text{为偶数}.
\end{cases}
\]
则对每个 \(m\ge 1\) 与 \(q\ge 0\)，有 \(\Omega_m(q)\in\ZZ\)，并满足统一的二阶线性递推
\[
\boxed{\;
\Omega_m(q)=L_m\,\Omega_m(q-1)-(-1)^m\,\Omega_m(q-2),\qquad
\Omega_m(0)=1,\ \Omega_m(1)=L_m\; }.
\]
尤其地，\(\Omega_1(q)=F_{q+1}\)、\(\Omega_2(q)=F_{2q+2}\) 为该机制的低阶退化。
\end{conclusion}
\begin{proof}
记 \(\widehat\varphi:=-\varphi^{-1}\) 为 \(\QQ(\sqrt5)\) 的共轭根，则
\(
L_m=\varphi^m+\widehat\varphi^{\,m}=\varphi^m+(-1)^m\varphi^{-m}
\)
且
\(
\varphi^m\widehat\varphi^{\,m}=(-1)^m
\)。
因此数列 \(q\mapsto \varphi^{mq}\) 与 \(q\mapsto \widehat\varphi^{\,mq}\) 均满足线性递推
\(
x_q=L_m x_{q-1}-(-1)^m x_{q-2}
\)；
由定义可将 \(\Omega_m(q)\) 写成二者的 \(\QQ\)-线性组合，从而递推式成立。
初值由 \(q=0,1\) 的直接代入得到。
递推系数与初值均为整数，故 \(\Omega_m(q)\in\ZZ\) 由归纳法给出。
最后，\(m=1,2\) 的退化式可由相同递推与初值对齐（或由 \(\Omega_1\)、\(\Omega_2\) 的闭式与 Binet 公式直接化简）得到。证毕。
\end{proof}

\begin{conclusion}[低阶迹矩的完全指数分解：\texorpdfstring{$\Tr((T^{(q)})^m)$}{Tr((T^(q))^m)} 的闭式]\label{con:xi-terminal-replica-softcore-low-moment-trace-exponential-decomposition}
对任意 \(q\ge 2\)，\(T^{(q)}\) 的低阶迹矩满足严格闭式
\[
\boxed{\;
\Tr\!\bigl((T^{(q)})^3\bigr)=\frac18\Big(8^{q}+3\cdot 6^{q}+3\cdot 5^{q}+4^{q}\Big)\; } ,
\]
\[
\boxed{\;
\Tr\!\bigl((T^{(q)})^4\bigr)=\frac1{16}\Big(16^{q}+4\cdot 12^{q}+4\cdot 10^{q}+2\cdot 9^{q}+4\cdot 8^{q}+7^{q}\Big)\; } ,
\]
\[
\boxed{\;
\Tr\!\bigl((T^{(q)})^5\bigr)=\frac1{32}\Big(32^{q}+5\cdot 24^{q}+5\cdot 20^{q}+5\cdot 18^{q}+5\cdot 16^{q}+5\cdot 15^{q}+5\cdot 13^{q}+11^{q}\Big)\; } .
\]
\end{conclusion}
\begin{proof}[证明要点]
由 \(T^{(q)}_{u,v}=\frac12\bigl(1+\mathbf{1}_{\{u\wedge v=0\}}\bigr)\) 展开
\(
\Tr((T^{(q)})^m)=\sum_{u_1,\dots,u_m}\prod_{k=1}^{m}T^{(q)}_{u_k,u_{k+1}}
\)
（其中 \(u_{m+1}=u_1\)），得到
\[
\Tr\!\bigl((T^{(q)})^m\bigr)
=2^{-m}\sum_{E'\subseteq E(C_m)}\ \sum_{u_1,\dots,u_m\in U_q}\ \prod_{(k,k+1)\in E'}\mathbf{1}_{\{u_k\wedge u_{k+1}=0\}}.
\]
对固定 \(E'\)，指示乘积在各坐标上分解，故内层求和按坐标完全因子化为 \(a(E')^q\)，其中 \(a(E')\) 为单坐标二元 \(m\)-环上满足相应“禁相邻 \(11\)”约束的可行赋值数。
当 \(m=3,4,5\) 时，\(E'\subseteq E(C_m)\) 的同构类型有限，可逐型枚举得到对应的 \(a(E')\) 分别为
\[
8,6,5,4;\qquad
16,12,10,9,8,7;\qquad
32,24,20,18,16,15,13,11,
\]
且其出现次数分别为
\[
1,3,3,1;\qquad
1,4,4,2,4,1;\qquad
1,5,5,5,5,5,5,1.
\]
代回并收集同底数项即得三条闭式。证毕。
\end{proof}

\begin{conclusion}[例外谱三次幂和的全闭式：\texorpdfstring{$\{8,6,5\}$}{\{8,6,5\}} 指数族与 Lucas--Pell 纠正项]\label{con:xi-terminal-replica-softcore-exceptional-power-sum-m3}
定义整数序列
\[
\Omega_3(q):=\frac{\varphi^{3q+6}+(-1)^q\varphi^{-3q}}{\varphi^{6}+1}\in\ZZ.
\]
则对一切 \(q\ge 2\) 有严格恒等式
\[
\boxed{\;
S_3(q)=\frac18\Big(8^{q}+3\cdot 6^{q}+3\cdot 5^{q}+\Omega_3(q)\Big)\; }.
\]
并且 \(\Omega_3\) 满足二阶线性递推
\[
\Omega_3(q)=4\,\Omega_3(q-1)+\Omega_3(q-2),\qquad \Omega_3(0)=1,\ \Omega_3(1)=4.
\]
\end{conclusion}
\begin{proof}[证明要点]
由定理 \ref{thm:pom-replica-softcore-characteristic-factorization}，\(T^{(q)}\) 的全谱由例外谱与固定谱簇并置给出。
因此
\(
\Tr((T^{(q)})^3)=S_3(q)+\sum_{j=0}^{q}(\binom{q}{j}-1)\bigl(\frac{(-1)^j\varphi^{q-2j}}{2}\bigr)^3
\)。
对固定谱簇幂和作二项式求和与几何级数求和可得
\[
\sum_{j=0}^{q}\Bigl(\binom{q}{j}-1\Bigr)\Bigl(\frac{(-1)^j\varphi^{q-2j}}{2}\Bigr)^3
=\frac18\bigl(4^q-\Omega_3(q)\bigr),
\]
其中 \(\varphi^3-\varphi^{-3}=4\) 被识别为 Lucas 数 \(L_3\)。
再与结论 \ref{con:xi-terminal-replica-softcore-low-moment-trace-exponential-decomposition} 的 \(\Tr((T^{(q)})^3)\) 闭式合并即得所示。
\(\Omega_3\) 的递推为结论 \ref{con:xi-terminal-replica-softcore-lucas-pell-omega-unified} 在 \(m=3\) 的特例。证毕。
\end{proof}

\begin{conclusion}[例外谱四次幂和的全闭式：\texorpdfstring{$\{16,12,10,9,8\}$}{\{16,12,10,9,8\}} 指数族与 Lucas--Pell 纠正项]\label{con:xi-terminal-replica-softcore-exceptional-power-sum-m4}
定义整数序列
\[
\Omega_4(q):=\frac{\varphi^{4q+8}-\varphi^{-4q}}{\varphi^{8}-1}\in\ZZ.
\]
则对一切 \(q\ge 2\) 有严格恒等式
\[
\boxed{\;
S_4(q)=\frac1{16}\Big(16^{q}+4\cdot 12^{q}+4\cdot 10^{q}+2\cdot 9^{q}+4\cdot 8^{q}+\Omega_4(q)\Big)\; }.
\]
并且 \(\Omega_4\) 满足二阶线性递推
\[
\Omega_4(q)=7\,\Omega_4(q-1)-\Omega_4(q-2),\qquad \Omega_4(0)=1,\ \Omega_4(1)=7.
\]
\end{conclusion}
\begin{proof}[证明要点]
同上，\(
\Tr((T^{(q)})^4)=S_4(q)+\sum_{j}(\binom{q}{j}-1)\bigl(\frac{(-1)^j\varphi^{q-2j}}{2}\bigr)^4
\)。
固定谱簇幂和化简为 \(\frac1{16}(7^q-\Omega_4(q))\)（其中 \(\varphi^4+\varphi^{-4}=7=L_4\)），再与结论 \ref{con:xi-terminal-replica-softcore-low-moment-trace-exponential-decomposition} 的 \(\Tr((T^{(q)})^4)\) 闭式合并即得所示。
递推为结论 \ref{con:xi-terminal-replica-softcore-lucas-pell-omega-unified} 在 \(m=4\) 的特例。证毕。
\end{proof}

\begin{conclusion}[例外谱五次幂和的全闭式：\texorpdfstring{$\{32,24,20,18,16,15,13\}$}{\{32,24,20,18,16,15,13\}} 指数族与 Lucas--Pell 纠正项]\label{con:xi-terminal-replica-softcore-exceptional-power-sum-m5}
定义整数序列
\[
\Omega_5(q):=\frac{\varphi^{5q+10}+(-1)^q\varphi^{-5q}}{\varphi^{10}+1}\in\ZZ.
\]
则对一切 \(q\ge 2\) 有严格恒等式
\[
\boxed{\;
S_5(q)=\frac1{32}\Big(32^{q}+5\cdot 24^{q}+5\cdot 20^{q}+5\cdot 18^{q}+5\cdot 16^{q}+5\cdot 15^{q}+5\cdot 13^{q}+\Omega_5(q)\Big)\; }.
\]
并且 \(\Omega_5\) 满足二阶线性递推
\[
\Omega_5(q)=11\,\Omega_5(q-1)+\Omega_5(q-2),\qquad \Omega_5(0)=1,\ \Omega_5(1)=11.
\]
\end{conclusion}
\begin{proof}[证明要点]
同理，固定谱簇五次幂和化简为 \(\frac1{32}(11^q-\Omega_5(q))\)（其中 \(\varphi^5-\varphi^{-5}=11=L_5\)），再与结论 \ref{con:xi-terminal-replica-softcore-low-moment-trace-exponential-decomposition} 的 \(\Tr((T^{(q)})^5)\) 闭式合并即可得到所示表达。
递推为结论 \ref{con:xi-terminal-replica-softcore-lucas-pell-omega-unified} 在 \(m=5\) 的特例。证毕。
\end{proof}

\endinput
